\documentclass[11pt]{article}
\usepackage[margin=1in]{geometry}
\usepackage{enumitem}
% =============================================================================
% Preambles/header.tex — shared LaTeX/Beamer preamble for this project
%
% Use from a Beamer lecture by prepending:
%     \input{header}
% and compiling with TEXINPUTS=../Preambles:$TEXINPUTS (the /compile-latex
% skill does this automatically).
%
% PALETTE CONTRACT. Color names defined here MUST match the names in
% Quarto/theme-template.scss so Beamer and Quarto renderings use the same
% palette. The scripts/check-palette-sync.sh script enforces this.
%
% To customize: edit the HEX values below AND the matching $variable in
% Quarto/theme-template.scss, then run scripts/check-palette-sync.sh.
% =============================================================================

% -----------------------------------------------------------------------------
% Engine detection + encoding
%
% Our /compile-latex skill uses XeLaTeX, where inputenc is unnecessary and
% can produce encoding warnings. Only load inputenc under pdfTeX.
% -----------------------------------------------------------------------------
\usepackage{iftex}
\ifPDFTeX
  \usepackage[utf8]{inputenc}
\fi

\usepackage{amsmath, amssymb, amsthm}
\usepackage{booktabs}
\usepackage{graphicx}

% -----------------------------------------------------------------------------
% PALETTE — mirrors Quarto/theme-template.scss variable names.
% Load xcolor and define colors BEFORE anything that references them
% (notably hyperref, hypersetup, and the Beamer color assignments).
% -----------------------------------------------------------------------------
\usepackage{xcolor}

\definecolor{primary-blue}{HTML}{012169}      % headings, accents
\definecolor{primary-gold}{HTML}{B9975B}      % emphasis, borders
\definecolor{highlight-yellow}{HTML}{F2A900}  % markers, alerts
\definecolor{light-bg}{HTML}{E8EDF5}          % title / section backgrounds
\definecolor{jet}{HTML}{1A1A1A}               % body text

% Semantic colors (used in diagrams and annotations)
\definecolor{positive}{HTML}{15803D}          % good / observed / identified
\definecolor{negative}{HTML}{B91C1C}          % bad / problematic / bias
\definecolor{neutral}{HTML}{525252}           % reference / context

% Highlight accents (match .hi-* classes in the SCSS)
\definecolor{hi-slate}{HTML}{314F4F}
\definecolor{hi-green}{HTML}{2E8B57}
\definecolor{hi-red}{HTML}{C41E3A}

% Semantic aliases so skill templates and snippets can write
% \textcolor{accent}{...} without hard-coding "primary-blue".
\colorlet{accent}{primary-blue}
\colorlet{accent2}{primary-gold}

% -----------------------------------------------------------------------------
% Hyperref — loaded AFTER xcolor + palette so link colors resolve.
% -----------------------------------------------------------------------------
\usepackage{hyperref}
\hypersetup{colorlinks=true, linkcolor=primary-blue, urlcolor=primary-blue,
            citecolor=primary-blue, pdfborder={0 0 0}}

% -----------------------------------------------------------------------------
% Course code — set once per deck (after \input{header}, before \begin{document})
% via \coursecode{CS401}. Shown in the footer on every slide so a deck's course
% is identifiable without opening the title page. Empty by default.
% -----------------------------------------------------------------------------
\makeatletter
\newcommand{\coursecode}[1]{\gdef\@coursecodetext{#1}}
\gdef\@coursecodetext{}
\makeatother

% -----------------------------------------------------------------------------
% Beamer theme assignments (only apply under Beamer).
%
% We detect Beamer via \@ifclassloaded{beamer}, which is the documented,
% non-internal way. This must be wrapped in \makeatletter/\makeatother
% because \@ifclassloaded contains an @-catcode character.
% -----------------------------------------------------------------------------
\makeatletter
\@ifclassloaded{beamer}{%
  \usefonttheme{professionalfonts}
  \setbeamercolor{structure}{fg=primary-blue}
  \setbeamercolor{frametitle}{fg=primary-blue}
  \setbeamercolor{title}{fg=primary-blue}
  \setbeamercolor{subtitle}{fg=primary-gold}
  \setbeamercolor{author}{fg=primary-blue}
  \setbeamercolor{institute}{fg=neutral}
  \setbeamercolor{date}{fg=primary-gold}
  \setbeamercolor{itemize item}{fg=primary-blue}
  \setbeamercolor{itemize subitem}{fg=primary-gold}
  \setbeamercolor{alerted text}{fg=primary-gold}
  \setbeamercolor{block title}{fg=white, bg=primary-blue}
  \setbeamercolor{block body}{bg=light-bg}
  % Semantic block variants: block=definition/concept (blue), exampleblock=
  % worked example (green/positive), alertblock=key takeaway or caution (gold).
  % Keep to <=2 colored boxes per slide across ALL three variants (INV-7).
  \setbeamercolor{block title example}{fg=white, bg=positive}
  \setbeamercolor{block body example}{bg=light-bg}
  \setbeamercolor{block title alerted}{fg=white, bg=primary-gold}
  \setbeamercolor{block body alerted}{bg=light-bg}

  % Minimal footer: course code (left) + page X / N (right), no navigation clutter
  \setbeamertemplate{navigation symbols}{}
  \setbeamertemplate{footline}{%
    \color{neutral}\footnotesize\hspace{0.7em}\@coursecodetext\hfill
    \insertframenumber/\inserttotalframenumber\hspace{0.7em}\vspace{0.3em}}
}{}
\makeatother

% -----------------------------------------------------------------------------
% TikZ setup — libraries the snippet gallery uses
% -----------------------------------------------------------------------------
\usepackage{tikz}
\usetikzlibrary{arrows.meta, positioning, calc, decorations.pathreplacing,
                fit, shapes.geometric, backgrounds}

% Shared TikZ styles — snippets and hand-written diagrams can pick these up
% via `\begin{tikzpicture}[every node/.style={...}]` or by naming specific
% styles (e.g., `\node[dag-node] (X) at (0,0) {X};`).
\tikzset{
  % Node styles for causal DAGs and similar
  dag-node/.style = {
    draw=primary-blue, fill=light-bg, rounded corners=2pt,
    minimum width=1.2cm, minimum height=0.9cm, align=center, font=\small
  },
  decision-node/.style = {
    draw=primary-gold, fill=white, diamond, aspect=1.8,
    minimum width=1.8cm, minimum height=1.2cm, align=center, font=\small,
    inner sep=1pt
  },
  % Edge styles — observed vs counterfactual
  observed-edge/.style = {-{Latex[length=2.5mm]}, thick, draw=jet},
  counterfactual-edge/.style = {-{Latex[length=2.5mm]}, thick, draw=neutral, dashed},
  confound-edge/.style = {-{Latex[length=2.5mm]}, thick, draw=negative, dashed},
  % Data-point styles for plots
  observed-dot/.style = {fill=primary-blue, draw=primary-blue, circle, inner sep=1.2pt},
  counterfactual-dot/.style = {fill=white, draw=neutral, circle, inner sep=1.2pt}
}

% -----------------------------------------------------------------------------
% Convenience macros
% -----------------------------------------------------------------------------
\newcommand{\muted}[1]{\textcolor{neutral}{#1}}
\newcommand{\key}[1]{\textcolor{primary-gold}{\textbf{#1}}}
\newcommand{\good}[1]{\textcolor{positive}{#1}}
\newcommand{\bad}[1]{\textcolor{negative}{#1}}

% Standout frame macro: full-bleed dark background for section transitions.
% Beamer-only (uses \begin{frame}/\setbeamercolor/\usebeamerfont) -- do not
% call from an article-class file (e.g. Notes/<CODE>/*-notes.tex). \muted,
% \key, \good, \bad, and \coursecode above are plain xcolor/gdef and are
% safe to use from either class; verified when Notes/article-class support
% was added.
% Expand: \transitionslide{Your transition title}
\newcommand{\transitionslide}[1]{%
  \begingroup
  \setbeamercolor{background canvas}{bg=primary-blue}
  \begin{frame}[plain]
    \centering
    \vfill
    {\usebeamerfont{title}\color{white}\Large #1\par}
    \vfill
  \end{frame}
  \endgroup
}

% Section divider frame (Beamer-only), an alternative to \transitionslide with a
% darker two-line style: near-black (jet) background, a small white label line
% above a larger gold title, and a thin gold rule along the bottom edge.
% Expand: \sectiondivider{Section 2}{Pointers}
\newcommand{\sectiondivider}[2]{%
  \begingroup
  \setbeamercolor{background canvas}{bg=jet}
  \begin{frame}[plain]
    \vspace*{3.0cm}
    {\color{white}\ttfamily\large #1\par}
    \vspace{0.5cm}
    {\color{primary-gold}\LARGE #2\par}
    \begin{tikzpicture}[remember picture, overlay]
      \draw[primary-gold, line width=2.5pt]
        ($(current page.south west)+(0,0.1)$) -- ($(current page.south east)+(0,0.1)$);
    \end{tikzpicture}
  \end{frame}
  \endgroup
}

% -----------------------------------------------------------------------------
% Channel watermark — "Concepts That Click" (YouTube).
%
% Article-class files (CompetitiveExam Q/A sets, Notes, Assignments) get a
% light diagonal draft watermark on every page plus a centered footer naming
% the channel. Beamer decks get a subtle TikZ background watermark instead
% (the footline already carries the course code). Centralized here so every
% MCQ PDF is branded without per-file edits.
% -----------------------------------------------------------------------------
\makeatletter
\@ifclassloaded{article}{%
  \usepackage{draftwatermark}
  \SetWatermarkText{Concepts That Click}
  \SetWatermarkScale{1}
  \SetWatermarkFontSize{2cm}
  \SetWatermarkLightness{0.86}
  \SetWatermarkAngle{45}
  \usepackage{fancyhdr}
  \pagestyle{fancy}
  \fancyhf{}
  \fancyfoot[C]{\footnotesize\textcolor{neutral}{Concepts That Click~|~YouTube: \href{https://www.youtube.com/@concepts.thatclick}{@concepts.thatclick}}}
  \renewcommand{\headrulewidth}{0pt}
  \renewcommand{\footrulewidth}{0pt}
  \fancypagestyle{plain}{%
    \fancyhf{}%
    \fancyfoot[C]{\footnotesize\textcolor{neutral}{Concepts That Click~|~YouTube: \href{https://www.youtube.com/@concepts.thatclick}{@concepts.thatclick}}}%
    \renewcommand{\headrulewidth}{0pt}%
    \renewcommand{\footrulewidth}{0pt}%
  }
}{}
\@ifclassloaded{beamer}{%
  \setbeamertemplate{background}{%
    \begin{tikzpicture}[remember picture, overlay]
      \node[rotate=45, opacity=0.055, align=center]
        at (current page.center)
        {\Huge\textcolor{neutral}{Concepts That Click}};
    \end{tikzpicture}%
  }
}{}
\makeatother 
\coursecode{JKSSB}
\renewcommand{\thesection}{53.\arabic{section}}
\newtheorem{example}{Example}[section]
\newenvironment{definitionbox}[1]{%
  \par\noindent\textbf{Definition (#1).}\ \itshape
}{\par\vspace{0.3em}}
\title{Networking Legacy Services and Web Standards --- Detailed Lecture Notes}
\author{JKSSB Computer Awareness}
\date{\today}

\begin{document}
\maketitle

\section{How to read this lesson}
This is a trap-drill lesson. Almost every concept here has a near-neighbour:
another term that sounds similar or sits next to it in a list but does a
different job. The exam builds its wrong options out of exactly those
neighbours, and it often asks the question with a negative stem (NOT, EXCEPT,
or an incorrect statement). The defence is always the same: state what the
term is, what it is not, and which neighbour it is confused with.

\begin{center}
\begin{tabular}{p{0.24\textwidth}p{0.66\textwidth}}
\toprule
\textbf{Term} & \textbf{Meaning in this lesson} \\
\midrule
Circuit switching & A communication method that reserves a dedicated path for the whole session before data moves. \\
Packet switching & A communication method that breaks data into packets, each travelling independently with no reserved path. \\
Telnet & TELecommunication NETwork; a remote-login service for running commands on a distant computer. \\
FTP & File Transfer Protocol; the service for uploading and downloading files. \\
SMTP & Simple Mail Transfer Protocol; sends mail from the sender toward the receiver's server. \\
POP3 & Post Office Protocol version 3; downloads mail to the local machine. \\
IMAP & Internet Message Access Protocol; keeps mail on the server and syncs it across devices. \\
W3C & World Wide Web Consortium; the body that develops open web standards. \\
DHTML & Dynamic HyperText Markup Language; HTML combined with styling and scripting for dynamic pages. \\
URL & Uniform Resource Locator; the full address of a web resource. \\
DNS & Domain Name System; translates domain names into IP addresses. \\
IP address & The numeric address of a device on the network. \\
OSI model & The seven-layer networking reference model. \\
Physical port & A hardware socket on the machine. \\
Logical port & A number identifying a service or process in networking software. \\
\bottomrule
\end{tabular}
\end{center}

\section{Packet switching vs circuit switching}
\textbf{Circuit switching} establishes a dedicated communication path between
the two ends before any data flows, and that path stays reserved for the whole
session. The classic example is the old telephone call: the line is yours
until you hang up, whether or not anyone is speaking.

\begin{definitionbox}{Circuit switching}
A switching method in which a dedicated path is set up and reserved for the
entire communication session before data transfer begins.
\end{definitionbox}

\textbf{Packet switching} works the opposite way. Data is broken into small
units called packets, and each packet travels independently through the
network. No path is reserved in advance; packets from many senders share the
same links. The Internet runs on packet switching.

\begin{definitionbox}{Packet switching}
A switching method in which data is divided into packets that travel
independently over shared links, with no dedicated path reserved in advance.
\end{definitionbox}

\begin{example}
An item states, ``The Internet reserves a dedicated circuit for each browsing
session.'' This is wrong. The Internet uses packet switching: pages arrive as
independent packets over shared links, and no path is reserved for the session.
The wrong option has described circuit switching instead.
\end{example}

\section{Telnet and FTP}
\textbf{Telnet} (TELecommunication NETwork) is a remote-login service. It lets
a user connect to another computer over the network and run commands there as
if sitting at its keyboard. Telnet sends its traffic, including passwords, in
plain text. It is not a file-transfer service and not a mail service.

\textbf{FTP} (File Transfer Protocol) is the file-transfer service. It moves
files between computers: uploading from the local machine to a server and
downloading back. It is not a remote-login shell and not a mail protocol.

The exam tests the triple together: Telnet logs in remotely, FTP moves files,
and SMTP sends mail. Any option that swaps these roles --- ``Telnet transfers
files'' or ``FTP sends mail'' --- is a planted near-neighbour distractor.

\begin{definitionbox}{Telnet}
A network service that provides remote login to another computer so the user
can run commands on it.
\end{definitionbox}

\begin{definitionbox}{FTP}
A network protocol used to transfer files between computers on a network.
\end{definitionbox}

\section{SMTP, POP3, and IMAP}
Mail handling is split across three protocols, and the exam tests the split.
\textbf{SMTP} (Simple Mail Transfer Protocol) \emph{sends} mail: it carries a
message from the sender's machine toward the receiver's mail server.
\textbf{POP3} (Post Office Protocol version 3) \emph{receives by download}:
it pulls mail down to the local machine, typically removing it from the
server. \textbf{IMAP} (Internet Message Access Protocol) \emph{receives by
sync}: it keeps mail on the server so the same mailbox can be read from
several devices.

\begin{definitionbox}{SMTP / POP3 / IMAP roles}
SMTP sends mail; POP3 downloads mail to the local machine; IMAP keeps mail on
the server and synchronises it across devices.
\end{definitionbox}

Three attached traps come from real papers. First, an option saying
``POP3 keeps mail synced on the server'' is wrong; that synced behaviour is
IMAP (TRAP-17). Second, ``a digital signature encrypts the body so only the
recipient can read it'' is wrong: a digital signature proves authenticity and
integrity, while encryption provides confidentiality (TRAP-15). Third,
``BCC hides the sender and the recipients from each other'' is wrong: BCC
hides recipients from each other, the sender always knows everyone, and To
and Cc names remain visible to all recipients (TRAP-16). Finally, in any mail
address the \textbf{@} sign separates the user name from the domain name
(TRAP-28).

\begin{example}
An item asks which protocol downloads mail to the local machine. The answer
is POP3. If the stem instead asks which keeps the mailbox synced on the
server across a phone and a desktop, the answer is IMAP, and if it asks
which sends the message in the first place, the answer is SMTP.
\end{example}

\section{W3C and DHTML}
The \textbf{W3C} (World Wide Web Consortium) develops the open standards of
the World Wide Web, such as HTML and CSS. It is a standards body --- not a
browser, not a search engine, and not an Internet service provider. The
near-neighbour is the \textbf{ISP} (Internet Service Provider), which sells
Internet access; W3C writes the standards that the access carries.

\begin{definitionbox}{W3C}
The World Wide Web Consortium: the international body that develops open
standards for the World Wide Web.
\end{definitionbox}

\textbf{HTML} (HyperText Markup Language) marks up the static structure of a
web page. \textbf{DHTML} (Dynamic HyperText Markup Language) is not a
separate replacement language; it is the combination of HTML with styling and
scripting that makes pages respond and change. The trap option calls DHTML a
brand-new language that replaces HTML. The correct statement is that DHTML
builds on HTML to add dynamic behaviour.

\section{URL, DNS, and IP addresses}
A \textbf{URL} (Uniform Resource Locator) is the full address of a resource
on the web. A \textbf{domain name} is the human-readable part inside it. The
\textbf{DNS} (Domain Name System) translates that human-readable name into a
numeric address. The \textbf{IP address} is that numeric address: it
identifies a device on the network.

The chain in one line is: type a URL, DNS resolves its domain name, and the
browser reaches the IP address. The exam swaps the links: an option saying
``DNS is the numeric address'' confuses DNS with the IP address, and one
saying ``the URL is the numeric address'' confuses the URL with the IP
address.

\section{IPv4 vs IPv6}
\textbf{IPv4} uses 32-bit addresses, written as four decimal groups. The
address space proved too small, so \textbf{IPv6} was introduced with 128-bit
addresses and a fixed header of 40 bytes (PYQ-18).

\begin{definitionbox}{IPv4 / IPv6 address size}
IPv4 addresses are 32 bits long; IPv6 addresses are 128 bits long.
\end{definitionbox}

The single most repeated distractor on this topic claims IPv6 is 256-bit
(TRAP-18, from JA 2026 Q78). It is always wrong: IPv6 is 128-bit, never
256-bit. A second attached trap from the same paper concerns companion
protocols: \textbf{ICMP} is a control and diagnostic protocol that carries no
application data and uses no ports (TRAP-19); \textbf{BGP} routes between
autonomous systems (inter-domain) and does not route inside one local network
(TRAP-20); and \textbf{HTTP/3} runs over QUIC over UDP, not over TCP
(TRAP-21, from JA 2026 Q79).

\section{The seven OSI layers}
The \textbf{OSI model} (Open Systems Interconnection) is the seven-layer
reference model for networking. From top to bottom the layers are
Application, Presentation, Session, Transport, Network, Data Link, and
Physical. A memory hook is ``All People Seem To Need Data Processing''.

\begin{definitionbox}{OSI model}
A seven-layer reference model for network communication: Application,
Presentation, Session, Transport, Network, Data Link, Physical.
\end{definitionbox}

The near-neighbour trap names four layers. Four is the layer count of the
TCP/IP model, not the OSI model. When the stem says OSI, the number is
always seven; an option offering four has imported the wrong model.

\section{Physical vs logical ports}
A \textbf{physical port} is a hardware socket on the machine: the place where
a cable or device plugs in. A \textbf{logical port} is a number in networking
software that identifies a service or process. The trap swaps them: a logical
port number is not a hardware socket, and a physical socket is not a service
number.

\begin{definitionbox}{Physical vs logical port}
A physical port is a hardware socket for cables or devices; a logical port is
a number identifying a network service or process.
\end{definitionbox}

\section{Exam retrieval and common confusions}
Nine distinctions prevent most errors on this topic.
\begin{enumerate}
  \item Circuit switching reserves a dedicated path; packet switching sends
    independent packets with no reserved path.
  \item Telnet gives remote login; FTP transfers files; SMTP sends mail. Do
    not swap the three roles.
  \item POP3 downloads mail to the local machine; IMAP keeps mail synced on
    the server (TRAP-17).
  \item A digital signature proves authenticity and integrity; encryption
    provides confidentiality (TRAP-15).
  \item BCC hides recipients from each other; the sender always knows all
    recipients (TRAP-16).
  \item IPv6 is 128-bit with a 40-byte header, never 256-bit (TRAP-18).
  \item ICMP carries no application data; BGP is inter-domain; HTTP/3 uses
    QUIC over UDP (TRAP-19, TRAP-20, TRAP-21).
  \item The OSI model has seven layers; four layers is the TCP/IP model.
  \item A physical port is a hardware socket; a logical port is a service
    number.
\end{enumerate}

\begin{example}
An item asks which statement is incorrect. Work each option against its
near-neighbour: does the switching option reserve or share the path, does
the protocol option match the service to its role, does the address option
use 128-bit for IPv6, does the layer option say seven for OSI? The option
that breaks its pair is the incorrect statement the stem asks for.
\end{example}

\subsection*{Exercises}
\begin{enumerate}[label=\arabic*.]
  \item Distinguish circuit switching from packet switching in one sentence
    each.
  \item Match Telnet, FTP, and SMTP to their roles.
  \item State the roles of SMTP, POP3, and IMAP, and say which one keeps mail
    synced on the server.
  \item Expand W3C, DHTML, URL, DNS, FTP, and SMTP.
  \item State the bit lengths of IPv4 and IPv6, and name the OSI layers from
    top to bottom.
\end{enumerate}

\subsection*{Solutions}
\begingroup\small
\begin{enumerate}[label=\arabic*.]
  \item Circuit switching reserves a dedicated path for the whole session
    before data moves; packet switching breaks data into independently
    travelling packets with no reserved path.
  \item Telnet gives remote login to another computer; FTP transfers files
    between computers; SMTP sends mail toward the receiver's server.
  \item SMTP sends mail, POP3 downloads it to the local machine, and IMAP
    keeps it on the server and syncs it; the synced-on-server protocol is
    IMAP.
  \item W3C --- World Wide Web Consortium; DHTML --- Dynamic HyperText Markup
    Language; URL --- Uniform Resource Locator; DNS --- Domain Name System;
    FTP --- File Transfer Protocol; SMTP --- Simple Mail Transfer Protocol.
  \item IPv4 is 32-bit and IPv6 is 128-bit; the OSI layers from the top are
    Application, Presentation, Session, Transport, Network, Data Link,
    Physical.
\end{enumerate}
\endgroup
\end{document}
